\section{Course Description}

This course examines the qualitative and quantitative \textbf{relationships} between the \textbf{rates and mechanisms} of organic reactions, and the \textbf{electronic and physical structures} of reactants. Among the topics considered are: theory and applications of inductive and resonance effects, \textbf{linear free energy relationships}, kinetic isotope effects, solvent effects, steric effects in substitution and elimination reactions, acids and bases and pericyclic reactions, applications of semi-empirical and \textit{ab initio} molecular orbital calculations.

\smallskip\noindent\textbf{Prerequisite}: Chemistry 2420, 3420

\section{Where \& When}

We are scheduled to meet at 8:30--9:20 in KCI 202 MWF.

\section{Instructor}

Dr. Barry Linkletter

\begin{flushleft}
BSc. -- University of Prince Edward Island\\
PhD. -- McGill University (Jik Chin)\\
Post Doc. -- University of California at Santa Barbara (T.C. Bruice)\\[1em]

Room 205\\
KCI Chemistry Building\\
628-4328\\
\texttt{blinkletter@upei.ca}
\end{flushleft}

\section{Office Hours}

I have an ``open door policy'' so please come by whenever I am available. However, if you want to be sure to meet at a particular time just contact me to make an appointment.

\newpage
\section{Textbook}

The textbook for this course is an \textbf{excellent book} and worth getting a copy and keeping it forever. There is a chapter for almost every field of chemistry in it. I will be using

\begin{quote}
\textbf{Modern Physical Organic Chemistry} by Eric V. Anslyn and Dennis A. Dougherty. University Science Books.\\
\href{https://mitpress.mit.edu/9781891389313/modern-physical-organic-chemistry/}{MIT Press Webpage}
\end{quote}




\begin{quote}
\textbf{Note}: The publisher of this textbook, \textit{University Science Books}, has recently been \textbf{purchased} by the \textit{MIT Press}. The website has \textbf{changed} and some links to the textbook in this webbook may be broken. Please \textbf{let me know} if you find a broken link.
\end{quote}

Access to the textbook material can be found in the links below or in the actual physical textbook.

\begin{itemize}
\item \textbf{UPEI Bookstore}:\\
\href{https://upei.bookware3000.ca/CourseSearch/?course[]=1,2025F,CHEM,CHEM441001,}{Order textbook from the UPEI Bookstore}

\item If that link does not work, you can purchase the \textbf{e-book} directly from\\
\href{https://www.vitalsource.com/en-ca/products/modern-physical-organic-chemistry-eric-v-anslyn-dennis-a-v9781891389306}{VitalSource}.

\item A copy is available \href{https://libraryupei.ca/chemistry-4410-physical-organic-chemistry-fall-2025}{on reserve at the library}.

\item A copy is available on the shelves in the \textbf{chemistry lounge}. Please leave it there \textbf{for all to use}.

\item \textbf{Front Matter of Textbook} (Moodle). This is the table of contents and the \textbf{preface} from the textbook. It was available on the former publisher's website. MIT press no longer provides it and I have made a copy available in the course webbook,

\item \textbf{Appendix 5 Problems} (Moodle). I have copied these \textbf{questions from the textbook} and made them available in the course webbook. Please obtain access to the textbook soon as I won't be able to do this again. Respect copyright---someday you may be an author.
\end{itemize}

\subsection{Other Textbook Suggestions}
The following books are also helpful:

\begin{itemize}
\item \textbf{McMurry's Organic Chemistry}\\
\href{https://openstax.org/details/books/organic-chemistry}{OpenStax Organic Chemistry}\\
A free textbook for 2\textsuperscript{nd}-year organic chemistry.

\item \textbf{Physical Organic Chemistry}, by Neil Isaacs. Available in the UPEI Library stacks: \textbf{QD476.I846}

\item \textbf{Organic Chemistry: An Intermediate Text}\\
\href{https://onlinelibrary-wiley-com.proxy.library.upei.ca/doi/book/10.1002/0471648736}{Wiley Online Library} via UPEI library online holdings\\
by Robert Hoffman. A good source for reviewing 3\textsuperscript{rd}-year advanced organic chemistry.

\item \textbf{Mechanism and Theory in Organic Chemistry}, 3\textsuperscript{rd} Edition, by Lowry and Richardson. Available in the UPEI Library stacks: \textbf{QD476.L68}

\item \textbf{Structure and Mechanism in Organic Chemistry}, by Felix A. Caroll

\item \textbf{Advanced Organic Chemistry: Reactions and Mechanisms}, by Bernard Miller; especially Chapters 1 and 5.

\item Oxford Chemistry Primers \#81: \textbf{Structure and Reactivity in Organic Chemistry}, by Howard Maskill.

\item Oxford Chemistry Primers \#45: \textbf{Mechanisms of Organic Reactions}, by Howard Maskill.
\end{itemize}

\section{Method of Delivery}

The course will be presented in a \textbf{topic format}. Most lesson topics will span three to six classes except where holidays or storms interfere. The first subject is review and will span the first five class meetings.

Course materials and communications are available through Moodle. The course will be organized in a topic format with three or more \textbf{in-person} class meetings.

\section{Evaluation}

Grading will follow the scheme outlined below. This grading scheme is open for revision. Any ideas?

\subsection{Grading Scheme}

\begin{center}
\begin{tabular}{ll}
\hline
Component & Grade \\
\hline
Class Participation & 10\% \\
Assignments & 40\% \\
Personal Exploration Report & 20\% \\
Mid-Term Tests & 30\% \\
\hline
\end{tabular}
\end{center}

Active \textbf{class participation} is required in this course. Students are expected to arrive prepared to participate in discussion of the topic. \textbf{Readings and exercises} will be assigned in advance.

The \textbf{assigned} textbook \textbf{questions} are meant to stimulate discussion and develop your \textbf{skills}. You may work together on problems, just thank those who helped you when presenting your solution. Take note that \textbf{the assignments are to be completed individually}.

Each week will feature an \textbf{assignment} or other goal to be completed. Most of these assignments/goals will involve writing a short \textbf{report}. Students will have at least 4 days to \textbf{complete} any assignment and \textbf{submit} the report.

Late reports may not be accepted unless \textbf{arrangements} are made in advance with a valid reason.

During the course, you will begin an exercise called the \textbf{personal exploration}. You will choose a \textbf{topic} that features physical organic chemistry and \textbf{explore} the current \textbf{literature}. This exercise will be \textbf{developed} over four weeks of the course and, in the end, you will produce a major \textbf{report} on your activities.

\subsection{Mid-Term Test Schedule}

\begin{center}
\begin{tabular}{ll}
\hline
Event & Date \\
\hline
Test \#1 & Oct. 26\textsuperscript{th} \\
Test \#2 & Nov. 23\textsuperscript{th} \\
\hline
\end{tabular}
\end{center}

There are \textbf{two mid-term tests} planned. They will be written during class time. There is no \textbf{final exam}.

I am open to adjustments of the grading scheme. Suggestions are welcome.

\section{Academic Integrity}

Academic integrity is important. As a \textbf{community of scholars}, the University of Prince Edward Island is committed to the principle of academic integrity among all its participants.

Academic \textbf{dishonesty} as defined in Academic Regulation 20 will not be tolerated.

The Department of Chemistry follows the principles of \textbf{progressive} discipline in response to concerns surrounding \textbf{academic integrity}.

For \textbf{reports and assignments}, the first incidence of academic misconduct will result in a verbal warning with a report to the Department Head. Students will have 3 days to resubmit the assignment for a grade up to 50\% of the original mark.

A \textbf{second incidence} of academic dishonesty will result in a grade of zero on the assignment with a report to the Dean of Science. A \textbf{third violation} will result in zero in the course.

All students must have a current \href{https://library.upei.ca/ai/home}{Academic Integrity Moodle Badge} to receive a grade. Nothing will be graded unless the \textbf{badge} is complete.

For \textbf{tests or exams}, the \textbf{first} incident of academic dishonesty will result in zero for the assessment and the \textbf{second} will result in zero for the course. Both will result in a report to the Dean of Science.

\section{Appeals}

If you feel that you have been treated unfairly in a course there is an avenue for appeal. The Department of Chemistry has an \textbf{academic appeals policy} that is available from the departmental office.

\section{Accommodation Statement}

Students may request \textbf{accommodation} as a result of \textbf{barriers} experienced related to disability, or any characteristic protected under Canadian Human Rights legislation.

Students who require academic accommodation for either classroom participation or the writing of tests and exams should make their request to \textbf{Accessibility Services} as soon as possible.

Please visit \href{https://upei.ca/accessibility}{UPEI Accessibility Services} for more information.

\section{Land Acknowledgement}

The land on which the University of Prince Edward Island is located is \textbf{Epekwitk}, the traditional and current territory of the \textbf{Mi'kmaq People} of this region. It is one of seven traditional Mi'kmaq regions that comprise Atlantic Canada, parts of eastern Quebec, and Maine.

Mi'kma'ki is the ancestral and unceded territory of the Mi'kmaq people, who, in 1725 first signed the \textbf{Treaties of Peace and Friendship} with the British Crown. Those treaties did not deal with the surrender of lands and resources, but instead recognized Mi'kmaq title and negotiated a path toward an ongoing relationship between nations.

We acknowledge that we carry out our daily work and learning in Mi'kma'ki.

\section{Storm Cancellations}

If UPEI is closed for inclement weather, no in-person class will be held that day. Should a storm cancellation fall on the date of an examination or assessment, the date will be moved to the next available class.

\section{Exam Absences}

If you miss a test or exam because of illness, death in the family, etc., please let the instructor know \textbf{as soon as possible} so that alternate arrangements can be made.

\section{Communications}

UPEI business is conducted by email. Your \textbf{UPEI email address is your official point of contact}. Please check your UPEI email often.

Emails will be responded to during standard working hours, 9\,am--5\,pm.

\textbf{Please} provide instructors, coordinators, teaching assistants, and staff with at minimum 48\,h to \textbf{respond} before sending a follow-up email.

